% !TeX program = xelatex
\documentclass[a4paper,10pt]{article}

% ============================================================================
% សន្លឹកកិច្ចការសិស្ស ស-ខ្មៅ (Student Worksheet - Grade 8 3.2.6 CINT)
% បរិបទ៖ គ្មានស្លាយ LCD និងគ្មានម៉ាស៊ីនគិតលេខ — សិស្សបំពេញផ្ទាល់លើសន្លឹកកិច្ចការ ២ ទំព័រ A4
% ============================================================================
\usepackage[a4paper, left=0.85cm, right=0.85cm, top=0.70cm, bottom=0.95cm, footskip=0.48cm]{geometry}
\usepackage{amsmath, amssymb}
\usepackage{unicode-math}
\usepackage{fontspec}
\usepackage{array, colortbl}
\usepackage[table]{xcolor}
\usepackage{enumitem}
\usepackage{pifont}
\usepackage{fancyhdr}
\usepackage{lastpage}
\usepackage{tikz}
\usepackage[most]{tcolorbox}
\usetikzlibrary{calc, arrows.meta, decorations.pathreplacing, positioning, babel}

\makeatletter
\let\@expl@@@mark@update@singlecol@structures@@\relax
\makeatother

\providecommand{\NGSFindFonts}{%
  \IfFileExists{C:/Windows/Fonts/cambria.ttc}{%
    \def\NGSSysPath{C:/Windows/Fonts/}%
    \def\NGSUserPath{C:/Users/KIT/AppData/Local/Microsoft/Windows/Fonts/}%
  }{%
    \IfFileExists{./fonts/KhmerOSsiemreap.ttf}{\def\NGSSysPath{./fonts/}\def\NGSUserPath{./fonts/}}{%
    \IfFileExists{../fonts/KhmerOSsiemreap.ttf}{\def\NGSSysPath{../fonts/}\def\NGSUserPath{../fonts/}}{%
    \IfFileExists{../../fonts/KhmerOSsiemreap.ttf}{\def\NGSSysPath{../../fonts/}\def\NGSUserPath{../../fonts/}}{%
    \IfFileExists{../../../fonts/KhmerOSsiemreap.ttf}{\def\NGSSysPath{../../../fonts/}\def\NGSUserPath{../../../fonts/}}{%
    \IfFileExists{../../../../fonts/KhmerOSsiemreap.ttf}{\def\NGSSysPath{../../../../fonts/}\def\NGSUserPath{../../../../fonts/}}{%
    \IfFileExists{../../../../../fonts/KhmerOSsiemreap.ttf}{\def\NGSSysPath{../../../../../fonts/}\def\NGSUserPath{../../../../../fonts/}}{%
    \IfFileExists{../../../../../../fonts/KhmerOSsiemreap.ttf}{\def\NGSSysPath{../../../../../../fonts/}\def\NGSUserPath{../../../../../../fonts/}}{%
      \def\NGSSysPath{../../../../../../../fonts/}\def\NGSUserPath{../../../../../../../fonts/}%
    }}}}}}}%
  }%
}
\NGSFindFonts

\setmathfont[Path={\NGSSysPath}, FontIndex=1, Scale=1.02]{cambria.ttc}
\defaultfontfeatures{Ligatures=TeX}

\setmainfont[
  Path={\NGSUserPath}, Extension=.ttf, Script=Khmer,
  BoldFont={KhmerOSsiemreap}, BoldFeatures={FakeBold=1.8},
  ItalicFont={KhmerOSsiemreap}, ItalicFeatures={FakeSlant=0.22}
]{KhmerOSsiemreap}

\newfontfamily\muol[Path={\NGSUserPath}, Extension=.ttf, Script=Khmer]{KhmerOSmuollight}
\newfontfamily\TimesEn[Path={\NGSSysPath}, Extension=.ttf, BoldFont={timesbd}, ItalicFont={timesi}]{times}

\XeTeXlinebreaklocale "km"
\XeTeXlinebreakskip = 0pt plus 1pt minus 0.1pt

\definecolor{panelgray}{gray}{0.93}
\definecolor{headgray}{gray}{0.85}

\pagestyle{fancy}
\fancyhf{}
\renewcommand{\headrulewidth}{0pt}
\fancyfoot[L]{\fontsize{7.5pt}{9.5pt}\selectfont អនុវិទ្យាល័យព្រះក្រសាំង \en{|} គណិតវិទ្យាថ្នាក់ទី៨ \en{|} មេរៀនទី៣៖ ២.៦ ការប្រាក់សមាស (គ្រូ៖ ហុក ប៊ុនថែន)}
\fancyfoot[R]{\fontsize{7.8pt}{9.5pt}\selectfont ទំព័រ \en{\thepage/\pageref{LastPage}}}

\setlength{\parindent}{0pt}
\setlength{\parskip}{0pt}
\setlength{\tabcolsep}{3.8pt}

\newcommand{\en}[1]{{\TimesEn #1}}
\newcommand{\enbf}[1]{{\TimesEn\bfseries #1}}
\newcommand{\unit}[1]{\,\text{\en{#1}}}
\newcommand{\cnum}[1]{\mbox{\fontsize{10pt}{12pt}\selectfont\ding{\numexpr181+#1\relax}}}
\newcommand{\chkbox}{\tikz[baseline=-0.12ex]\draw[black, line width=0.8pt] (0,0) rectangle (0.24cm,0.24cm);}
\newcommand{\blank}[1][1.5cm]{\makebox[#1]{\dotfill}}
\newcommand{\ansbox}[1][0.85cm]{\,\tikz[baseline=-0.5ex]\draw[black, line width=0.65pt, fill=white] (0,0) rectangle (#1, 0.42cm);\,}
\newcommand{\secbar}[2]{%
  \begin{tcolorbox}[enhanced, sharp corners, colback=headgray, colframe=black, boxrule=0.9pt, boxsep=1pt, top=2pt, bottom=2pt, left=5pt, right=5pt, before skip=4pt, after skip=3pt]
    {\muol\fontsize{9.2pt}{13pt}\selectfont #1}\hfill{\fontsize{8pt}{11pt}\selectfont\textbf{#2}}
  \end{tcolorbox}%
}

\newcolumntype{L}[1]{>{\raggedright\arraybackslash}p{#1}}
\newcolumntype{C}[1]{>{\centering\arraybackslash}p{#1}}

\begin{document}
\fontsize{9.6pt}{14.4pt}\selectfont

% ============================================================================
% ទំព័រទី ១៖ E1 + E2 (POE 1 & POE 2) + E3 សំយោគលំហូរ និងធនាគារហេតុផល
% ============================================================================
\begin{tcolorbox}[enhanced, sharp corners, colback=white, colframe=black, boxrule=1.0pt, boxsep=2pt, top=3pt, bottom=3pt, left=6pt, right=6pt]
  \begin{center}
    {\muol\fontsize{10.0pt}{15pt}\selectfont អនុវិទ្យាល័យព្រះក្រសាំង \enbf{--} សន្លឹកកិច្ចការសិស្ស \enbf{(}មេរៀនទី៣៖ ២\enbf{.}៦ ការប្រាក់សមាស\enbf{)}}\\[3pt]
    {\fontsize{8.8pt}{12pt}\selectfont
    ឈ្មោះសិស្ស / ដៃគូ៖ \blank[5.2cm] \hspace{0.25cm}
    ថ្នាក់ទី៨៖ \blank[1.1cm] \hspace{0.25cm}
    ក្រុមទី៖ \blank[1.0cm] \hspace{0.25cm}
    កាលបរិច្ឆេទ៖ ០២/១២/២០២៦}
  \end{center}
\end{tcolorbox}

% --- E1 + E2 ---
\secbar{ផ្នែកទី១៖ បេសកកម្មទស្សន៍ទាយ \enbf{(E1)} និងការរុករកជាដៃគូ \enbf{(E2 -- POE 1 \& POE 2)}}{រយៈពេល ១៤ នាទី}

\begin{tcolorbox}[enhanced, sharp corners, colback=panelgray, colframe=black, boxrule=0.75pt, boxsep=1.5pt, top=2.5pt, bottom=2.5pt, left=5pt, right=5pt]
  \textbf{បេសកកម្ម \en{E1} (គូសធីក \chkbox\ មុនពិសោធន៍)៖}
  (១) មីងសៅដាក់ប្រាក់ $8\,000\,000\text{៛}$ អត្រាការប្រាក់សមាស $15\%$/ឆ្នាំ ($2$ ឆ្នាំ) — ការប្រាក់សរុបស្មើ៖ \chkbox\ $2\,400\,000\text{៛}$ \quad \chkbox\ $2\,580\,000\text{៛}$\\
  (២) ប្រាក់លើស $+180\,000\text{៛}$ នៅឆ្នាំទី ២ កើតចេញពី៖ \chkbox\ ធនាគារថែមរង្វាន់ \quad \chkbox\ ការប្រាក់ $15\%$ គិតលើការប្រាក់ឆ្នាំទី ១
\end{tcolorbox}

\vspace{3pt}
\textbf{ក. សកម្មភាពទី ១ (\en{POE 1})៖ តារាងបង្វិលប្រាក់ដើមរបស់មីងសៅ ($P_1 = 8\,000\,000\text{៛}, R = 15\%$, ទំ.៣៧)}

{\renewcommand{\arraystretch}{1.48}
\begin{tabular}{|C{2.40cm}|C{3.60cm}|L{6.80cm}|C{4.80cm}|}
  \hline
  \rowcolor{headgray}
  \textbf{គ្រាប្រចាំឆ្នាំ} & \textbf{ប្រាក់ដើមដើមគ្រា} & \textbf{ការប្រាក់ប្រចាំឆ្នាំ ($I_k = \frac{P_k \times 15 \times 1}{100}$)} & \textbf{ប្រាក់សរុបចុងគ្រា ($P_{k+1}$)} \\
  \hline
  ឆ្នាំទី ១ & $P_1 = 8\,000\,000\text{៛}$ & $I_1 = \dfrac{8\,000\,000\times 15}{100} = \blank[1.7cm]\text{៛}$ & $P_2 = 8\,000\,000 + I_1 = \blank[1.5cm]$ \\
  \hline
  ឆ្នាំទី ២ & $P_2 = \blank[1.8cm]\text{៛}$ & $I_2 = \dfrac{\blank[1.7cm]\times 15}{100} = \blank[1.7cm]\text{៛}$ & $A_2 = P_2 + I_2 = \blank[1.8cm]\text{៛}$ \\
  \hline
  សរុប $2$ ឆ្នាំ & ធៀបសាមញ្ញ $2.4\text{លាន}$ & $I_{\text{សមាស}} = I_1 + I_2 = \blank[1.9cm]\text{៛}$ & សមាសលើស $+\blank[1.5cm]\text{៛}$ \\
  \hline
\end{tabular}}

\vspace{4pt}
$\bullet$ \textbf{ពន្យល់ (\en{E})៖} ប្រាក់លើស $+180\,000\text{៛}$ គឺជាការប្រាក់ $15\%$ គិតលើការប្រាក់ឆ្នាំទី ១ ($I_1 = \blank[1.8cm]\text{៛}$)!

\vspace{5pt}
\textbf{ខ. សកម្មភាពទី ២ (\en{POE 2})៖ បញ្ជីបង្វិលប្រាក់ដើម ៣ ឆ្នាំរបស់ទ្រាចាន់ ($P_1 = 3\,000\,000\text{៛}, R = 20\%$, គំរូ ១ ទំ.៣៧–៣៨)}

{\renewcommand{\arraystretch}{1.44}
\begin{tabular}{|C{2.40cm}|C{3.80cm}|L{6.60cm}|C{4.80cm}|}
  \hline
  \rowcolor{headgray}
  \textbf{គ្រាប្រចាំឆ្នាំ} & \textbf{ប្រាក់ដើមដើមគ្រា ($P_k$)} & \textbf{ការប្រាក់ប្រចាំឆ្នាំ ($20\%$ នៃ $P_k$)} & \textbf{ប្រាក់សរុបចុងឆ្នាំ ($P_k + I_k$)} \\
  \hline
  ឆ្នាំទី ១ & $3\,000\,000\text{៛}$ & $I_1 = 3\,000\,000 \times \frac{20}{100} = \blank[1.6cm]\text{៛}$ & $P_2 = 3\,600\,000\text{៛}$ \\
  \hline
  ឆ្នាំទី ២ & $3\,600\,000\text{៛}$ & $I_2 = 3\,600\,000 \times \frac{20}{100} = \blank[1.6cm]\text{៛}$ & $P_3 = \blank[2.0cm]\text{៛}$ \\
  \hline
  ឆ្នាំទី ៣ & $\blank[1.9cm]\text{៛}$ & $I_3 = \blank[1.7cm] \times \frac{20}{100} = \blank[1.5cm]\text{៛}$ & $A_3 = \blank[2.0cm]\text{៛}$ \\
  \hline
\end{tabular}}

\vspace{4pt}
$\bullet$ \textbf{ពន្យល់ (\en{E})៖} ការប្រាក់សមាស ៣ ឆ្នាំរបស់ទ្រាចាន់ស្មើ $I = I_1 + I_2 + I_3 = \blank[2.0cm]\text{៛}$ (រីឯការប្រាក់សាមញ្ញបានត្រឹម $1\,800\,000\text{៛}$)!

% --- E3 ---
\secbar{ផ្នែកទី២៖ សំយោគវិធីគណនា \enbf{(E3)} និងធនាគារហេតុផល \enbf{B\&W} \cnum{1}\enbf{--}\cnum{5}}{សៀវភៅពុម្ព ទំ.៣៧–៣៨}

\begin{tabular}{|L{9.20cm}|L{9.20cm}|}
  \hline
  \textbf{១. លំហូរគណនាការប្រាក់សមាសតាមគ្រា៖}\par\vspace{2.5pt}
  $\bullet$ \textbf{ការប្រាក់ក្នុងគ្រាទី $k$៖} $I_k = \dfrac{P_k \cdot R \cdot T_{\text{គ្រា}}}{100}$\par\vspace{3pt}
  $\bullet$ \textbf{ប្រាក់ដើមថ្មីសម្រាប់គ្រាបន្ទាប់៖} $P_{k+1} = P_k + \blank[0.9cm]$\par\vspace{3pt}
  $\bullet$ \textbf{ការប្រាក់សមាសសរុប៖} $I = I_1 + I_2 + \dots + I_n = A_n - \blank[0.8cm]$\par\vspace{3pt}
  $\bullet$ \textbf{គិតរាល់កន្លះឆ្នាំម្ដង ($6$ ខែ)៖} $T_{\text{គ្រា}} = \blank[0.8cm]$ ឆ្នាំ ($1$ ឆ្នាំ $= 2$ គ្រា)
  &
  \textbf{២. ធនាគារហេតុផល \en{B\&W} (\cnum{1}--\cnum{5})៖}\par\vspace{2.5pt}
  \cnum{1} បង្វិលប្រាក់ដើម៖ ចុងគ្រានីមួយៗ បូក $I_k$ ចូល $P_k$ ជាប្រាក់ដើមថ្មី ។\par\vspace{2.5pt}
  \cnum{2} គ្រាទី ១ ស្មើគ្នា៖ នៅគ្រាទី ១ សមាស និងសាមញ្ញស្មើគ្នាជានិច្ច ។\par\vspace{2.5pt}
  \cnum{3} ចាប់ពីគ្រាទី ២៖ សមាសធំជាងសាមញ្ញជានិច្ច (ការប្រាក់បង្កើតការប្រាក់) ។\par\vspace{2.5pt}
  \cnum{4} ការប្រាក់សរុប៖ $I = I_1 + I_2 + \dots + I_n$ ឬ $I = A_n - P_1$ ។\par\vspace{2.5pt}
  \cnum{5} គិតរាល់ $6$ ខែម្ដង៖ ប្រើ $T = \frac{1}{2}$ ឆ្នាំក្នុង ១ គ្រា ($1$ ឆ្នាំមាន $2$ គ្រា) ។ \\
  \hline
\end{tabular}

\vspace{4pt}
\begin{tcolorbox}[enhanced, sharp corners, colback=panelgray, colframe=black, boxrule=0.65pt, boxsep=1.5pt, top=2.5pt, bottom=2.5pt, left=5pt, right=5pt]
  \textbf{ប្រយ័ត្នអន្ទាក់៖} ហាមយកការប្រាក់គ្រាទី ១ ($I_1$) ទៅគុណនឹងចំនួនគ្រា! ត្រូវយក $P_2 = P_1 + I_1$ មកគិតការប្រាក់គ្រាទី ២ ជានិច្ច!
\end{tcolorbox}

\vspace{4pt}
\textbf{៣. លំហាត់ត្រួតពិនិត្យរហ័ស៖ គណនាប្រាក់ដើមបង្វិល និងការប្រាក់សមាស ២ គ្រា}

\vspace{2.5pt}
{\renewcommand{\arraystretch}{1.60}
\begin{tabular}{|C{4.40cm}|C{4.40cm}|C{4.40cm}|C{4.40cm}|}
  \hline
  \rowcolor{panelgray}
  $P_1=100\$, 10\%$៖ $P_2=\blank[0.7cm]$ & $P_2=110\$, 10\%$៖ $I_2=\blank[0.7cm]$ & $P_1=100\$, 10\%, 2\text{ឆ្នាំ}\to I=\blank[0.6cm]$ & សមាសលើសសាមញ្ញ៖ $+\blank[0.6cm]\$$ \\
  \hline
  $P_1=200\$, 20\%$៖ $P_2=\blank[0.7cm]$ & $P_2=240\$, 20\%$៖ $I_2=\blank[0.7cm]$ & $P_1=200\$, 20\%, 2\text{ឆ្នាំ}\to A_2=\blank[0.6cm]$ & $P_1=1\,000\$, 10\%, 2\text{ឆ្នាំ}\to I=\blank[0.5cm]$ \\
  \hline
  គិតរាល់ $6$ ខែ៖ $T=\blank[0.7cm]$ ឆ្នាំ & ក្នុង $1$ ឆ្នាំ ($6$ខែម្ដង) មាន $\blank[0.5cm]$ គ្រា & ក្នុង $2$ ឆ្នាំ ($6$ខែម្ដង) មាន $\blank[0.5cm]$ គ្រា & $P_1=500\$, 10\%, 2\text{ឆ្នាំ}\to I=\blank[0.5cm]$ \\
  \hline
  $P_1=400\$, 10\%, 2\text{ឆ្នាំ}\to I=\blank[0.5cm]$ & $P_1=1\,000\$, 20\%\to P_2=\blank[0.6cm]$ & $P_2=1\,200\$, 20\%\to I_2=\blank[0.5cm]$ & $P_1=800\$, 10\%, 6\text{ខែ}\to I_1=\blank[0.5cm]$ \\
  \hline
\end{tabular}}

\newpage

% ============================================================================
% ទំព័រទី ២៖ E4 លំហាត់គំរូ + E5 ប្រតិបត្តិ + វិនិច្ឆ័យ A-D + Exit Check + កិច្ចការផ្ទះ
% ============================================================================
\secbar{ផ្នែកទី៣៖ ពង្រីកចំណេះដឹងលើការប្រាក់សមាសគិតរាល់កន្លះឆ្នាំម្ដង \enbf{(E4 --} គំរូ ២ ទំ\enbf{.}៣៨\enbf{)}}{រយៈពេល ១០ នាទី}

\textbf{លំហាត់គំរូទី ២ (សៀវភៅពុម្ព ទំ.៣៨ — គិតការប្រាក់សមាសរាល់កន្លះឆ្នាំម្ដង)៖} គេយកប្រាក់ដើម $\mathbf{7\,500\,000\text{៛}}$ ទៅដាក់ធនាគារតាមអត្រាការប្រាក់សមាស $\mathbf{16\%}$ ក្នុងមួយឆ្នាំ រយៈពេល $\mathbf{1}$ \textbf{ឆ្នាំ} ដោយធនាគារបូកការប្រាក់បញ្ចូលប្រាក់ដើមរាល់ \textbf{កន្លះឆ្នាំម្ដង ($6$ ខែម្ដង)} ។ រកប្រាក់ការសមាសសរុប និងប្រាក់សរុបចុងឆ្នាំ៖

{\renewcommand{\arraystretch}{1.48}
\begin{tabular}{|L{4.50cm}|L{9.10cm}|C{4.20cm}|}
  \hline
  \rowcolor{headgray}
  \textbf{គ្រាបង្វិលកន្លះឆ្នាំ} & \textbf{ការគណនាប្រាក់ការ និងប្រាក់ដើមថ្មីដោយដៃ} & \textbf{លទ្ធផល និងហេតុផល} \\
  \hline
  \textbf{១. កំណត់ចំនួនគ្រា និង $T$} & $6\text{ ខែ} = \blank[0.8cm]$ ឆ្នាំ $\implies$ ក្នុង $1$ ឆ្នាំ មានចំនួន \blank[0.7cm] គ្រា & ច្រកកន្លះឆ្នាំ (\cnum{5}) \\
  \hline
  \textbf{២. គ្រាទី ១ ($6$ ខែដំបូង)} ($P_1 = 7\,500\,000\text{៛}$) & $I_1 = \dfrac{7\,500\,000\times 16\times \frac{1}{2}}{100} = \blank[1.5cm]\text{៛} \Rightarrow P_2 = \blank[1.7cm]\text{៛}$ & បង្វិលដើមថ្មី $P_2$ (\cnum{1}) \\
  \hline
  \textbf{៣. គ្រាទី ២ ($6$ ខែបន្ទាប់)} ($P_2 = \blank[1.7cm]\text{៛}$) & $I_2 = \dfrac{\blank[1.7cm]\times 16\times \frac{1}{2}}{100} = \blank[1.6cm]\text{៛}$ & ការប្រាក់គ្រាទី ២ (\cnum{3}) \\
  \hline
  \textbf{៤. ប្រាក់ការ និងប្រាក់សរុប} & $I = I_1 + I_2 = \blank[1.7cm]\text{៛}$ \;;\; $A_2 = P_2 + I_2 = \blank[1.7cm]\text{៛}$ & សរុបចុងឆ្នាំ (\cnum{4}) \\
  \hline
\end{tabular}}

\vspace{4pt}
$\bullet$ \textbf{សំណួរឆ្លុះបញ្ចាំង \en{(EIA Looking-Back)}៖} បើផ្ញើប្រាក់សន្សំ គួរជ្រើសការប្រាក់ \blank[1.6cm] ; បើខ្ចីប្រាក់គេ គួរជ្រើសការប្រាក់ \blank[1.6cm] ។

% --- E5 ---
\secbar{ផ្នែកទី៤៖ វាយតម្លៃលទ្ធផលសិក្សា \enbf{(E5 --} ប្រតិបត្តិ ទំ\enbf{.}៣៨ \enbf{+} វិនិច្ឆ័យ \enbf{A--D +} សំបុត្រចេញ\enbf{)}}{រយៈពេល ១០ នាទី}

\textbf{ក. លំហាត់ប្រតិបត្តិ (សៀវភៅពុម្ព ទំ.៣៨ និងពង្រឹងប្រៀបធៀប — ធ្វើជាដៃគូ រួចប្ដូរសន្លឹកផ្ទៀង \en{Pause-Review})៖}

{\renewcommand{\arraystretch}{1.48}
\begin{tabular}{|L{8.90cm}|L{8.90cm}|}
  \hline
  \textbf{១) ប្រតិបត្តិ (ទំ.៣៨ — ប្រាក់បញ្ញើសមាស ២ ឆ្នាំ)៖}\par\vspace{2pt}
  គេយកប្រាក់ $\mathbf{2\,700\,000\text{៛}}$ ទៅដាក់ធនាគារតាមអត្រាការប្រាក់សមាស $\mathbf{20\%}$ ក្នុងមួយឆ្នាំ រយៈពេល $\mathbf{2}$ ឆ្នាំ ។ គណនាប្រាក់ការសមាស និងប្រាក់សរុប៖\par\vspace{2pt}
  $\bullet$ ឆ្នាំទី ១៖ $I_1 = \dfrac{2\,700\,000\times 20}{100} = \blank[1.5cm]\text{៛} \Rightarrow P_2 = \blank[1.6cm]$\par\vspace{2pt}
  $\bullet$ ឆ្នាំទី ២៖ $I_2 = \dfrac{\blank[1.5cm]\times 20}{100} = \blank[1.5cm]\text{៛}$\par\vspace{2pt}
  $\bullet$ សរុប៖ $I = \blank[1.8cm]\text{៛}$ \;;\; $A_2 = \blank[1.9cm]\text{៛}$ (\cnum{4})
  &
  \textbf{២) ពង្រឹង (ប្រៀបធៀបការប្រាក់សមាស ទល់នឹង សាមញ្ញ)៖}\par\vspace{2pt}
  មីងណារីដាក់ប្រាក់ $\mathbf{5\,000\,000\text{៛}}$ តាមអត្រា $\mathbf{10\%}$ ក្នុងមួយឆ្នាំ រយៈពេល $\mathbf{2}$ ឆ្នាំ ។ ប្រៀបធៀបប្រាក់ការទាំងពីរប្រភេទ៖\par\vspace{2pt}
  $\bullet$ \textbf{ក. ការប្រាក់សាមញ្ញ៖} $I_{\text{សាមញ្ញ}} = \dfrac{5\,000\,000\times 10\times 2}{100} = \blank[1.5cm]$\par\vspace{2pt}
  $\bullet$ \textbf{ខ. ការប្រាក់សមាស៖} $I_1 = 500\,000\text{៛}, P_2 = \blank[1.6cm]\text{៛}$ $\Rightarrow I_2 = \blank[1.4cm]\text{៛} \Rightarrow I_{\text{សមាស}} = \blank[1.6cm]\text{៛}$ (លើស $+\blank[1.1cm]$) \\
  \hline
\end{tabular}}

\vspace{5pt}
\textbf{ខ. គូសវង់ជុំវិញជម្រើសត្រឹមត្រូវតែ ១ គត់ ក្នុងចំណោម \en{A, B, C, D} រួចសរសេរលេខហេតុផល \cnum{1}--\cnum{5} នៃជម្រើសខុស៖}

{\renewcommand{\arraystretch}{1.36}
\begin{tabular}{|C{4.40cm}|C{4.40cm}|C{4.40cm}|C{4.40cm}|}
  \hline
  \cellcolor{panelgray}\textbf{ជម្រើស \en{A}} & \cellcolor{panelgray}\textbf{ជម្រើស \en{B}} & \cellcolor{panelgray}\textbf{ជម្រើស \en{C}} & \cellcolor{panelgray}\textbf{ជម្រើស \en{D}} \\
  \hline
  $P_1=2.7\text{លាន}, 20\%, 2$ ឆ្នាំ $\Rightarrow I = 1\,188\,000\text{៛}$ & $P_1=2.7\text{លាន}, 20\%, 2$ ឆ្នាំ $\Rightarrow I = 1\,080\,000\text{៛}$ & គិតសមាសរាល់ $6$ ខែម្ដង ជំនួស $T = 6$ ក្នុងមួយគ្រា & នៅឆ្នាំទី ១ ការប្រាក់សមាសធំជាងសាមញ្ញ \\[3pt]
  \chkbox\ ត្រូវ \quad \chkbox\ ខុស (\blank[0.7cm]) & \chkbox\ ត្រូវ \quad \chkbox\ ខុស (\blank[0.7cm]) & \chkbox\ ត្រូវ \quad \chkbox\ ខុស (\blank[0.7cm]) & \chkbox\ ត្រូវ \quad \chkbox\ ខុស (\blank[0.7cm]) \\
  \hline
\end{tabular}}

\vspace{5pt}
\begin{tcolorbox}[enhanced, sharp corners, colback=white, colframe=black, boxrule=0.85pt, boxsep=2.5pt, top=3.5pt, bottom=3.5pt, left=5pt, right=5pt]
  {\muol\fontsize{8.8pt}{12.5pt}\selectfont សំបុត្រចេញរហ័ស \enbf{(Exit Check --} ៦០ វិនាទី\enbf{)}៖}\\[2pt]
  ដាក់ប្រាក់ $1\,000\,000\text{៛}$ អត្រាការប្រាក់សមាស $10\%$/ឆ្នាំ រយៈពេល $2$ ឆ្នាំ៖ (១) $I_{\text{សមាស}} = \blank[1.7cm]\text{៛}$ \quad\en{|}\quad (២) $A_2 = \blank[1.8cm]\text{៛}$
\end{tcolorbox}

\vspace{3pt}
\begin{tcolorbox}[enhanced, sharp corners, colback=panelgray, colframe=black, boxrule=0.75pt, boxsep=2pt, top=3pt, bottom=3pt, left=5pt, right=5pt]
  {\muol\fontsize{8.8pt}{12.5pt}\selectfont កិច្ចការផ្ទះ \enbf{(}សៀវភៅគោល ទំព័រ ៤២ លំហាត់ \enbf{\#17} ក\enbf{,} ខ និង \enbf{\#18)}៖}\\[2pt]
  $\bullet$ \textbf{លំហាត់ \#17 (ក, ខ — ទំ.៤២)៖} គណនាការប្រាក់សមាស និងប្រាក់សរុបតាមតារាងបង្វិលប្រាក់ដើម ។\\
  $\bullet$ \textbf{លំហាត់ \#18 (ទំ.៤២)៖} គណនាការប្រាក់សមាសគិតរាល់កន្លះឆ្នាំម្ដង និងអានមុន \textbf{២.៧. ការទិញបណ្ដាក់ (ទំ.៣៨–៤០)} ។
\end{tcolorbox}

\end{document}
